\subsection{Experimental Results: Per-Thread Scheduling Strategies}
\label{sec:scheduling-results}

We evaluate the five scheduling strategies on two workloads using the Yahoo Streaming Benchmark (YSB) with 50 million tuples (4.8\,GB) and 10,000 distinct campaign IDs.
All experiments run on a dual-socket server with 2 NUMA nodes (32 logical CPUs each, 16 physical cores with hyperthreading) and 512\,GB RAM.
For thread counts $\leq 32$, processes are pinned to NUMA node~0 via \texttt{numactl --cpunodebind=0 --membind=0}.
At 64 threads, both NUMA nodes are used without pinning.

\subsubsection{Stateless Workload (Filter Only)}

The stateless workload applies a predicate (\texttt{WHERE event\_type < 1}) with approximately 80\% selectivity.
This is a single-pipeline query that produces \emph{no successor tasks}---each source buffer is processed in a single pipeline execution through filter and checksum sink.
The only scheduling-relevant contention is on the admission queue (source $\rightarrow$ workers).

Table~\ref{tab:stateless-strategies} shows the throughput in MB/s for each strategy.

\begin{table}[h]
\centering
\caption{Stateless workload throughput (MB/s) by strategy and thread count. Bold values indicate the best strategy per thread count.}
\label{tab:stateless-strategies}
\begin{tabular}{l|rrrrrrr}
\hline
\textbf{Strategy} & \textbf{1T} & \textbf{2T} & \textbf{4T} & \textbf{8T} & \textbf{16T} & \textbf{32T} & \textbf{64T} \\
\hline
\texttt{GLOBAL\_QUEUE}      & 72 & 134 & 241 & 403 & \textbf{564} & \textbf{593} & \textbf{388} \\
\texttt{PT\_ROUND\_ROBIN}   & 72 & 134 & 234 & 400 & 395 & 506 & 191 \\
\texttt{PT\_SMALLEST\_QUEUE}& 72 & 131 & 228 & 368 & 412 & 498 & 362 \\
\texttt{PT\_CHOOSE\_TWO}    & 72 & 133 & 238 & 400 & 400 & 496 & 355 \\
\texttt{PRODUCER\_LOCAL}    & 72 & \textbf{135} & \textbf{241} & \textbf{403} & 406 & 477 & 278 \\
\hline
\textit{Speedup vs.\ 1T}   & 1.0$\times$ & 1.9$\times$ & 3.3$\times$ & 5.6$\times$ & 7.8$\times$ & 8.2$\times$ & 5.4$\times$ \\
\hline
\end{tabular}
\end{table}

Key observations:

\begin{itemize}
    \item \textbf{At 1--8 threads, all strategies perform identically} (within 3\%), with \texttt{PRODUCER\_LOCAL} and \texttt{GLOBAL\_QUEUE} slightly ahead.
    At low thread counts, queue contention is negligible, so the dispatch policy has minimal impact.

    \item \textbf{\texttt{GLOBAL\_QUEUE} dominates at 16--64 threads}, achieving 564\,MB/s at 16T (7.8$\times$) and peaking at 593\,MB/s at 32T (8.2$\times$).
    Since this workload produces no successor tasks, the per-thread queue strategies add dispatch overhead without any locality benefit.
    The single shared queue provides perfect load balancing with no dispatch cost for internal tasks (there are none).

    \item \textbf{Per-thread strategies diverge at 16T and above.}
    \texttt{ROUND\_ROBIN} drops to 395\,MB/s at 16T and collapses to 191\,MB/s at 64T.
    The round-robin counter becomes a contention point when many threads increment it simultaneously, and the even distribution can leave some queues overfull while others are empty if thread execution speeds vary.
    \texttt{SMALLEST\_QUEUE} and \texttt{CHOOSE\_TWO} fare better (400--412\,MB/s at 16T) because they adapt to load imbalance, but still trail \texttt{GLOBAL\_QUEUE} due to the overhead of maintaining per-thread queues for a workload that does not benefit from them.

    \item \textbf{\texttt{PRODUCER\_LOCAL} offers no advantage over \texttt{GLOBAL\_QUEUE}} on this workload.
    Since there are no successor tasks, the producer-local policy is never triggered---all tasks come from the admission queue.
    The 477\,MB/s at 32T (vs.\ 593 for \texttt{GLOBAL\_QUEUE}) reflects the overhead of the per-thread queue infrastructure without any benefit.

    \item \textbf{The 32T $\rightarrow$ 64T throughput drop} is caused by two factors:
    (1)~at 32T, all threads share 16 physical cores via hyperthreading, reducing per-thread throughput;
    (2)~at 64T, NUMA pinning is disabled, introducing cross-node memory access latency.
    \texttt{GLOBAL\_QUEUE} handles this best (388\,MB/s) because its single shared queue naturally load-balances across both NUMA nodes.
\end{itemize}

\subsubsection{Stateful Workload (Aggregation with Tumbling Window)}

The stateful workload performs a grouped aggregation with a 30-second tumbling window (\texttt{GROUP BY campaign\_id WINDOW TUMBLING(current\_ms, SIZE 30 SEC)}).
This involves hash table lookups for 10,000 distinct groups and window state management.
The 30-second window produces approximately 1,667 window emissions across 50 million tuples, resulting in very few successor tasks (one emission per $\sim$30,000 input tuples).

Table~\ref{tab:stateful-strategies} shows the results.

\begin{table}[h]
\centering
\caption{Stateful workload throughput (MB/s) by strategy and thread count. Bold values indicate the best strategy per thread count.}
\label{tab:stateful-strategies}
\begin{tabular}{l|rrrrrrr}
\hline
\textbf{Strategy} & \textbf{1T} & \textbf{2T} & \textbf{4T} & \textbf{8T} & \textbf{16T} & \textbf{32T} & \textbf{64T} \\
\hline
\texttt{GLOBAL\_QUEUE}      & 77 & \textbf{112} & \textbf{142} & \textbf{147} & 133 & \textbf{179} & \textbf{143} \\
\texttt{PT\_ROUND\_ROBIN}   & 77 & 108 & 141 & 141 & \textbf{139} & 170 & 126 \\
\texttt{PT\_SMALLEST\_QUEUE}& 77 & 109 & 127 & 133 & 137 & 152 & 130 \\
\texttt{PT\_CHOOSE\_TWO}    & 78 & 110 & 135 & 144 & 138 & 171 & 122 \\
\texttt{PRODUCER\_LOCAL}    & 77 & 110 & 135 & 142 & 137 & 167 & 124 \\
\hline
\textit{Speedup vs.\ 1T}   & 1.0$\times$ & 1.4$\times$ & 1.8$\times$ & 1.9$\times$ & 1.8$\times$ & 2.3$\times$ & 1.8$\times$ \\
\hline
\end{tabular}
\end{table}

Key observations:

\begin{itemize}
    \item \textbf{Scaling is limited to 2.3$\times$}, with \texttt{GLOBAL\_QUEUE} achieving the peak of 179\,MB/s at 32T.
    The aggregation operator's hash table and window state management are the bottleneck, not the scheduling strategy.
    Adding more threads provides diminishing returns because the per-tuple cost of hash table lookups and state updates dominates.

    \item \textbf{All strategies perform within 10--15\% of each other} at every thread count.
    This confirms that for compute-bound stateful queries, the choice of scheduling strategy is secondary to operator efficiency.

    \item \textbf{\texttt{GLOBAL\_QUEUE} wins at most thread counts} (2T, 4T, 8T, 32T, 64T), consistent with the stateless results.
    Its advantage is most pronounced at 32T (179 vs.\ 152--171 for per-thread strategies), where the single shared queue avoids the per-thread overhead that provides no benefit for this workload.

    \item \textbf{The throughput dip at 16T} (133--139\,MB/s, below 8T at 141--147\,MB/s) is caused by increased contention on the shared aggregation hash table.
    With 16 threads all performing concurrent hash lookups and window state updates, synchronization overhead temporarily exceeds the parallelism gain.
    At 32T, throughput recovers as hyperthreading allows more threads to overlap memory stalls with computation.

    \item \textbf{\texttt{SMALLEST\_QUEUE} performs worst} at low thread counts (127\,MB/s at 4T vs.\ 142 for \texttt{GLOBAL\_QUEUE}).
    The $O(N)$ linear scan of all queue sizes adds measurable overhead when tasks are short (stateful tasks spend most time in the aggregation operator, leaving little room for scheduling overhead to hide).

    \item \textbf{\texttt{PRODUCER\_LOCAL} provides no measurable benefit} despite the stateful query producing window emission successor tasks.
    The reason is that window emissions are too rare (1,667 across 50M tuples) to create meaningful cache locality advantages.
    A workload with higher-frequency successor tasks (e.g., smaller window sizes or multi-stage pipelines) would be needed to observe the benefit.
\end{itemize}

\subsubsection{Summary}

\begin{itemize}
    \item \textbf{\texttt{GLOBAL\_QUEUE} is the best default strategy.}
    It wins or ties on both workloads at most thread counts.
    Its simplicity (no batch size tuning, no flags) and natural load balancing make it robust across workload types.

    \item \textbf{Per-thread queue strategies add overhead without benefit for single-pipeline queries.}
    The dispatch cost (atomic counter for round-robin, linear scan for smallest-queue, random sampling for choose-two) and per-thread queue maintenance are only justified when successor tasks are frequent enough to benefit from locality.

    \item \textbf{\texttt{PRODUCER\_LOCAL} is the best per-thread strategy} at low thread counts (1--8T), tying with \texttt{GLOBAL\_QUEUE}.
    However, it degrades at higher thread counts due to load imbalance when all admission tasks are distributed evenly but successor tasks accumulate on specific threads.

    \item \textbf{The scheduling strategy matters most at high thread counts} ($\geq 16$T), where contention on shared data structures amplifies the differences between strategies.
    At 1--8 threads, all strategies are within 5\% of each other.

    \item \textbf{Both workloads are bottlenecked by factors other than scheduling:}
    the stateless workload by admission queue throughput (source buffer production rate), and the stateful workload by operator execution time (hash table and window state management).
    To fully exercise the scheduling strategies, a multi-pipeline workload with frequent successor tasks is needed.
\end{itemize}
